\begin{claim}{C3.6}{Known}{Gauduchon metrics with specified Chern--Ricci form}
There exists a Gauduchon metric on $X$.  More precisely, for every
Gauduchon metric $\omega_0$ there is another Gauduchon metric
$\omega$ with $[\omega^2]_A=[\omega_0^2]_A$ and
\begin{equation}\label{eq:C-Ricci}
\operatorname{Ric}^{\mathrm{Ch}}(\omega)=\pi f^*\omega_{FS}.
\end{equation}
This Ricci form is semipositive and has rank one at regular points
of $f$.  No Chern--Ricci-flat Hermitian metric exists on $X$.
\end{claim}
\begin{proof}
Gauduchon's conformal theorem
\citep[Theorem 2, pp. 132--133]{Gauduchon1985} applies to the compact
connected complex manifold $X$.  In complex dimension three, it
produces a positive conformal multiple $\omega_0$ of any Hermitian
form with $\partial\bar\partial\omega_0^2=0$.
For clarity, its reduction is a scalar elliptic equation: if
$L\phi$ is a nonzero fixed normalization of
$\operatorname{tr}_\omega(i\partial\bar\partial\phi)$, then
$L$ has kernel the constants by the maximum principle and Fredholm
index zero.  The one-dimensional kernel of $L^*$ has a strictly
positive generator $u$.  Indeed a sign-changing generator would
admit a nonzero nonnegative function orthogonal to it, placing such
a function in $\im L$, contrary to the maximum principle; the strong
maximum principle then gives strict positivity.  The form
$u^{1/2}\omega$ is Gauduchon.

The canonical-bundle formula in Part A and $4S_2=f^*(p_2)$ give
$K_X^2=f^*\OO(-1)$.  Consequently
$c_{1,BC}(T_X)=\tfrac12[\theta]$ with
$\theta=f^*\omega_{FS}$.  We use the conventional unnormalized
Chern--Ricci form $-i\partial\bar\partial\log\det g$,
which represents $2\pi c_{1,BC}(T_X)$.  Thus the real closed form
$\pi\theta$ represents its class.  Applying
\citet[Theorem 1.3]{STW2017} with background $\omega_0$ and prescribed
form $\pi\theta$ proves \eqref{eq:C-Ricci} and preservation of the
Aeppli class.  Finally C3.1 shows $[\theta]_{BC}\ne0$.  A metric
with zero Chern--Ricci form would make $c_{1,BC}(T_X)=0$, which is
impossible.
\end{proof}
